\section{Method}
\label{sec:method}

We describe TTMG as three coupled components: the claim schema and the
graph that stores them (\S\ref{ssec:schema}), the write-time linker that
induces truth-maintenance edges (\S\ref{ssec:linker}), and the
read-time truth-consistent retriever (\S\ref{ssec:retrieve}).
Figure~\ref{fig:overview} sketches the data flow.

% Placeholder tikz overview figure. Rendered inline by main.tex.
\begin{figure}[t]
  \centering
  \begin{tikzpicture}[
    node distance=1.1cm,
    box/.style={draw,rounded corners,align=center,minimum width=2.6cm,minimum height=0.9cm,font=\small},
    memory/.style={draw,rounded corners,align=center,minimum width=3.4cm,minimum height=1.2cm,font=\small,fill=blue!5},
    link/.style={->,>=latex,thick},
    softlink/.style={->,>=latex,dashed},
  ]
    \node[box] (turn) {\textbf{Turn}\\(\emph{session}, \emph{speaker}, \emph{text})};
    \node[box,right=1.1cm of turn] (writer) {\textbf{Writer}\\claim extraction};
    \node[memory,right=1.1cm of writer] (claim) {\textbf{Claim}\\$(s,p,o,I,\pi,v,\rho,\phi)$};
    \node[memory,below=0.9cm of claim,fill=green!5] (graph) {\textbf{Memory graph}\\support / contradict / supersede};
    \node[box,left=1.1cm of graph] (linker) {\textbf{Linker}\\joint LLM labeller};

    \draw[link] (turn) -- (writer);
    \draw[link] (writer) -- (claim);
    \draw[link] (claim) -- (graph);
    \draw[link] (linker) -- (graph);
    \draw[softlink] (claim) -- (linker);

    \node[box,below=0.9cm of graph,fill=orange!5] (retrieve) {\textbf{Truth-consistent retrieval}\\query parse $\to$ $k$-NN $\to$ max-consistent subgraph};
    \node[box,below=0.5cm of retrieve] (reader) {\textbf{Reader}\\(answer or abstain)};
    \draw[link] (graph) -- (retrieve);
    \draw[link] (retrieve) -- (reader);
  \end{tikzpicture}
  \caption{TTMG at a glance. At write time, the \emph{writer} turns each
    session into \emph{claim} nodes with validity intervals and provenance;
    the \emph{linker} labels candidate neighbours with one of
    \texttt{support}, \texttt{contradict}, \texttt{supersede}, or
    \texttt{unrelated}. At read time, a query parser extracts the temporal
    anchor, $k$-NN retrieves candidates from the \emph{active} subset, and
    the retriever returns the maximum consistent subgraph. The reader
    receives claims with their validity intervals and can be forced to
    abstain when two active claims disagree.}
  \label{fig:overview}
\end{figure}


\subsection{Claim-level memory schema}
\label{ssec:schema}

The atomic unit of TTMG memory is a \emph{claim} $c = (t, s, p, o, I,
\pi, v, \rho, \phi)$, where $t$ is a natural-language paraphrase of the
statement; $(s, p, o)$ is an optional shallow subject--predicate--object
triple; $I = [\tau_{\text{from}}, \tau_{\text{to}}]$ is a validity
interval over the conversation's reference timeline (either side may
be open); $\pi \in \{\texttt{assert}, \texttt{deny}\}$ is the claim's
polarity; $v \in \{\texttt{stable}, \texttt{preference},
\texttt{state}\}$ is a coarse volatility class; $\rho$ is the
provenance record (session id, turn id, speaker, original timestamp);
and $\phi \in [0, 1]$ is the writer's confidence.

A turn of conversation is converted to zero or more claims by a single
LLM call that we call the \emph{writer}, \texttt{WriteClaims($T$)},
taking a whole session $T$ as input so that (i) the cost is one call
per session rather than per turn, and (ii) cross-turn coreference can
be resolved before extraction. The writer is prompted to skip small
talk, to default $\tau_{\text{from}}$ to the session timestamp, and to
tag corrections as \texttt{preference} claims with fresh
$\tau_{\text{from}}$.

Each claim $c$ is also associated with a normalised embedding
$\mathbf{e}(c) \in \mathbb{R}^{d}$ produced by a frozen bi-encoder. We
use a single index across all claims; active and superseded claims
coexist in the index but are filtered at query time.

\subsection{Write-time linker}
\label{ssec:linker}

When a new claim $c^{\star}$ is inserted, TTMG asks whether
$c^{\star}$ should be joined to any existing claim by one of four
labelled edges:
\begin{itemize}
  \item \texttt{support}: $c^{\star}$ and the neighbour co-assert a
  fact; both remain active.
  \item \texttt{contradict}: $c^{\star}$ and the neighbour disagree on
  the same $(s,p)$ at overlapping $I$; neither supersedes the other.
  \item \texttt{supersede}: $c^{\star}$ replaces the neighbour because
  it is a later or explicitly corrective assertion about the same
  $(s,p)$; the neighbour's \texttt{active} bit is cleared, and it is
  retained only for history-style queries.
  \item \texttt{unrelated}: the default; edges with this label are
  not materialised.
\end{itemize}

The linker proceeds in two stages. First, a cheap gate selects
candidate neighbours: we retain a neighbour only if (i) its cosine
similarity to $c^{\star}$ exceeds $\tau_{\text{sim}}$ \emph{and} (ii)
its subject or predicate tokens overlap with those of $c^{\star}$, or
cosine similarity alone exceeds a higher bypass threshold
$\tau_{\text{bypass}}$. This keeps the LLM linker off the path for
the overwhelming majority of unrelated new claims.

Second, the surviving candidates are joined into a single LLM call
that jointly labels all pairs and returns a label and a confidence for
each. If the linker returns \texttt{contradict} but the new claim is
strictly later on the same $(s,p)$, we upgrade to \texttt{supersede}
to enforce temporal monotonicity. Edges with confidence $\geq
\tau_{\text{hard}}$ are \emph{hard} edges that change retrieval
behaviour; below that threshold they are \emph{soft} edges that are
logged but have no effect on active-set filtering.

\subsection{Truth-consistent retrieval}
\label{ssec:retrieve}

Given a question $q$, the retriever performs four steps.

\textbf{1. Parse.} A small LLM call extracts $q$'s temporal anchor
(e.g., ``currently'', ``in 2024''), its entity, its intent
(\texttt{lookup}, \texttt{update\_check},
\texttt{temporal\_reasoning}, \texttt{small\_talk}), and two boolean
flags: $\texttt{asks\_history}$ (does the question ask about a past
state?) and $\texttt{asks\_current}$ (does it ask about the current
state?). These flags decide whether to retrieve only \texttt{active}
claims or to include superseded ancestors.

\textbf{2. k-NN.} We retrieve the top-$k$ claims by cosine similarity
to $q$, restricted to the active set when $\texttt{asks\_history} =
\texttt{false}$ and expanded to include superseded ancestors
otherwise. No cross-encoder reranker is required.

\textbf{3. Maximum-consistent subgraph.}
Given candidate claims $C = \{c_1, \dots, c_k\}$ with confidences
$\phi_i$ and the induced hard-\texttt{contradict} subgraph
$G_{\neg} = (C, E_{\neg})$, the retained set $R$ is the output of
Algorithm~\ref{alg:mcs}.

\begin{algorithm}[H]
  \caption{Greedy maximum-consistent subgraph}
  \label{alg:mcs}
  \begin{algorithmic}[1]
    \Require candidates $C$, contradict edges $E_{\neg}$, $k_{\text{keep}}$
    \State sort $C$ by $\phi_i$ descending, then by $\tau_{\text{from},i}$ descending
    \State $R \gets \emptyset$
    \For{$c \in C$}
      \If{$|R| \geq k_{\text{keep}}$} \State \textbf{break} \EndIf
      \If{$\forall c' \in R, (c, c') \notin E_{\neg}$} \State $R \gets R \cup \{c\}$ \EndIf
    \EndFor
    \State \Return $R$
  \end{algorithmic}
\end{algorithm}

\begin{definition}[Consistency of $R$]
  A set $R \subseteq C$ is \emph{consistent} if no two claims in $R$
  are joined by a hard-\texttt{contradict} edge: $\forall c, c' \in R:
  (c, c') \notin E_{\neg}$.
\end{definition}

\begin{proposition}
  The set $R$ returned by Algorithm~\ref{alg:mcs} is consistent. Its
  weight $\sum_{c \in R} \phi_c$ is within a factor $1/\Delta$ of the
  optimum, where $\Delta$ is the maximum degree of $G_{\neg}$
  restricted to $C$. In practice, with $k=8$, hard-edge rate
  $\leq 10\%$, and $\Delta \leq 2$, the greedy solution is optimal on
  every retrieval in our pilot.
\end{proposition}

Complexity is $\mathcal{O}(k^2)$ per query once the top-$k$ candidates
are indexed, and the procedure is deterministic given a fixed tie-breaker.

\textbf{4. Evidence-bound abstention.} If, among the top
$k_{\text{keep}}+2$ candidates, there exist two active claims sharing
an $(s,p)$ that are connected by a hard \texttt{contradict} edge and
neither supersedes the other, the retriever sets an
\texttt{abstain} flag. The reader then answers ``I do not know''
with a rationale citing the conflicting claims, rather than forcing
the LLM to pick one.

\paragraph{Hybrid reader context.}
Alongside the claim graph, TTMG maintains a frozen-embedding $k$-NN
store over \emph{raw turns}, built during ingestion with no extra LLM
calls. The reader receives both the truth-consistent claim subgraph
(preferred for temporal questions; carries validity intervals and
\texttt{active}/\texttt{past} flags) and the top-$k_{\text{raw}}$
raw turns (similarity-only fallback when claims are insufficient).
This preserves the TR advantage without losing the recall that flat
retrieval provides on general factual slices. Ablations in
\S\ref{sec:ablation} disable the linker to isolate its contribution.
